\documentclass[conference]{IEEEtran}

\usepackage{cite}
\usepackage{amsmath,amssymb,amsfonts}
\usepackage{algorithmic}
\usepackage{graphicx}
\usepackage{textcomp}
\usepackage{xcolor}
\usepackage{booktabs}
\usepackage{tagging}
\usepackage{hyperref}

\usepackage{sc26repro}

% Comment both lines for the final submission.
%\usetag{explanation}
%\usetag{example}

\begin{document}

\twocolumn[%
{\begin{center}
\Huge
Appendix: Artifact Description/Artifact Evaluation
\end{center}}
]

%%%%%%%%%%%%%%%%%%%%%%%%%%%%%%%%%%%%%%%%%%%%%%%%%%%%%
%  AD Appendix
%%%%%%%%%%%%%%%%%%%%%%%%%%%%%%%%%%%%%%%%%%%%%%%%%%%%%

\appendixAD

\section{Overview of Contributions and Artifacts}

\subsection{Paper's Main Contributions}

\begin{description}
\item[$C_1$] A theoretical proof, via blocked CDAGs and the Red-Blue-White Pebble Game, that a poor data layout can inflate the I/O lower bound by up to a factor of the cache-line length~$B$ relative to a layout-aware optimum, and that no schedule can close this gap.
\item[$C_2$] A composable layout algebra of five primitives (pad, permute, block, shuffle, zip) whose compositions suffice to express any bijective array rearrangement.
\item[$C_3$] Two heuristic cost metrics---the average new-block count~$\mu$ (guiding permutation) and the average block distance~$\Delta$ (guiding blocking)---together with a cookbook (Table~II) that maps observable performance symptoms to corrective primitives.
\item[$C_4$] Empirical validation on microbenchmarks and two production HPC codes---Quantum ESPRESSO and the ICON dynamical core (2025 ACM Gordon Bell Prize for Weather and Climate Modelling winner)---demonstrating up to $1.18\times$ speedup on GPUs and $1.12\times$ on CPUs.
\end{description}

\subsection{Computational Artifacts}

\begin{description}
\item[$A_1$] \url{https://doi.org/XX.XXXX/zenodo.XXXXXXX}
%% TODO: Replace with actual Zenodo DOI after archival
\end{description}

The artifact comprises two Git repositories and the DaCe framework:
\begin{itemize}
\item \textbf{R1:} \url{https://github.com/ThrudPrimrose/SC26-Layout-Artifacts} --- microbenchmarks, Quantum ESPRESSO case study, and the ICON single-loopnest benchmarks.
\item \textbf{R2:} \url{https://github.com/spcl/icon-artifacts} (branch \texttt{sc26}) --- ICON velocity-tendencies source SDFGs and the full-module lowering pipeline. A curated snapshot (SDFGs, transformation utilities, and driver scripts) is also mirrored into~R1 under \texttt{E6\_VelocityTendencies/full\_velocity\_tendencies/} for self-contained reproduction.
\item \textbf{R3:} \url{https://github.com/spcl/dace} (branches \texttt{yakup/dev} and \texttt{f2dace/staging}) --- the DaCe parallel programming framework with the layout-transformation infrastructure.
\end{itemize}

\begin{center}
\small
\begin{tabular}{clll}
\toprule
Exp. & Contributions & Paper Elements & Repo \\
\midrule
E1 & $C_3$, $C_4$         & Fig.~4             & R1 \\
E2 & $C_{2\text{--}4}$    & Fig.~8             & R1 \\
E3 & $C_{2\text{--}4}$    & Fig.~9, Tab.~III   & R1 \\
E4 & $C_{2\text{--}4}$    & Fig.~10            & R1 \\
E5 & $C_{2\text{--}4}$    & Fig.~11, Lst.~1    & R1 \\
E6 & $C_3$, $C_4$         & Figs.~12--13, Tab.~IV & R1, R2 \\
\bottomrule
\end{tabular}
\end{center}

All experiments depend on R3 (DaCe): E1--E5 and E6 single-loopnest benchmarks (loopnests~1--6) use branch \texttt{yakup/dev}; E6's full-module evaluation uses branch \texttt{f2dace/staging}.
Contribution~$C_1$ is a pen-and-paper proof with no computational artifact; the illustrative quotient-graph figures (Figures~2--3) are produced by a helper script in the supplemental material folder of~R1.

%%%%%%%%%%%%%%%%%%%%%%%%%%%%%%%%%%%%%%%%%%%%%%%%%%%%%%%%
\section{Artifact Identification}
%%%%%%%%%%%%%%%%%%%%%%%%%%%%%%%%%%%%%%%%%%%%%%%%%%%%%%%%

\newartifact

\artrel

Artifact~$A_1$ validates the layout algebra~($C_2$), cost metrics~($C_3$), and the cookbook-driven optimization workflow~($C_4$) through six experiments.
Experiments E1--E4 are microbenchmarks that isolate individual layout primitives and measure their effect on achieved memory bandwidth.
Experiment~E5 evaluates composed layout transformations on a production kernel from Quantum ESPRESSO.
Experiment~E6 applies the full optimization workflow to the ICON dynamical core, encompassing loop-nest canonicalization, per-pattern layout selection, cross-nest conflict resolution, and end-to-end evaluation of the velocity-tendencies module.

\artexp

\textbf{E1 (Matrix Addition, Figure~4):}
Elementwise addition $C \mathrel{+}= A + B$ with $B$ stored in column-major layout.
On our platforms, schedule-only optimization recovers 36\% (Zen4) to 72\% (Grace) of STREAM peak on CPUs and 78\% (MI300A) to 95\% (Hopper) on GPUs; applying layout transformations closes the gap to 68\% (Zen4) and 88\% (Grace) on CPUs and 99\% (both GPUs).
More broadly, on any platform with explicit scratchpad memory, schedule-only tiling into shared memory should recover most of the peak bandwidth.
On CPUs, where the programmer has no direct cache control, a persistent gap should remain that only layout transformations can close; its magnitude depends on hardware characteristics.

\textbf{E2 (Complex Conjugation, Figure~8):}
In-place conjugation of $P$ complex arrays under AoS, SoA, and AoSoA-16 layouts.
SoA should outperform AoS at low~$P$ due to fewer wasted cache-line bytes.
As $P$ grows, SoA may collapse on CPUs with limited L1D associativity because of set conflicts from page-aligned allocations; we observe this on Zen4 at $P \geq 15$ but not on Grace.
On both GPUs, SoA and AoSoA-16 sustain near-peak bandwidth, while AoS suffers from strided access to non-contiguous fields.
AoSoA-16 remains robust across all four backends and all~$P$ values, demonstrating the value of partial zipping when the optimal granularity cannot be selected statically.

\textbf{E3 (Matrix Transpose, Figure~9, Table~III):}
$N \times N$ fp32 transpose under row-major and blocked layouts.
On CPUs, row-major writes scatter at stride~$N$, exceeding prefetcher and TLB capacity; blocking improves bandwidth by $1.7\times$ on Zen4 and $3.2\times$ on Grace over row-major.
Our sweep reaches 61\% (Zen4) and 81\% (Grace) of STREAM peak, compared with HPTT's 44\% on both CPUs.
On GPUs, shared-memory tiling absorbs the stride: our sweep reaches 60\% (MI300A) and 85\% (Hopper) under row-major, matching vendor libraries; padding or shuffling remain necessary to avoid bank conflicts.
Table~III reports the reduction in~$\Delta$ from blocking at each cache-line size.
On any CPU without an explicit scratchpad, blocking should be the primary recourse for stride-heavy access patterns; on GPUs, the layout penalty is absorbed by shared-memory staging.

\textbf{E4 (Gather-Accumulate-Scatter, Figure~10):}
Indirect scatter--accumulate pattern $y[\sigma(i)] \mathrel{+}= x[i]$ before and after applying a shuffle (index sorting), evaluated under uniform-random and exact BaTiO$_3$ index distributions.
On MI300A GPU, the shuffle improves bandwidth from 21\% to 30\% (uniform) and from 23\% to 42\% (BaTiO$_3$) of STREAM peak.
On Hopper GPU, from 35\% to 62\% (uniform) and from 48\% to 82\% (BaTiO$_3$).
On CPUs, the improvement is modest across both distributions---9--15\% on Zen4 and 7--13\% on Grace---because the shuffle introduces a new indirect gather that limits prefetcher effectiveness.
The GPU gains should generalize to any coalescing-sensitive architecture; CPU gains depend on the trade-off between improved scatter locality and degraded gather locality.

\textbf{E5 (Quantum ESPRESSO \texttt{addusxx\_g}, Figure~11):}
Composed unzip, permute, and shuffle on the \texttt{addusxx\_g} kernel.
The composed transformations yield $1.12\times$ on Zen4 CPU, $1.11\times$ on Grace CPU, $1.18\times$ on Hopper GPU, and $1.00\times$ on MI300A GPU.
The improvement should generalize to any platform where the baseline Fortran AoS layout wastes cache-line bandwidth on the untouched real part (unzip), misaligns the innermost stride (permute), and scatters indirect writes (shuffle).
The MI300A GPU shows no gain because its unified memory architecture already absorbs the access pattern efficiently.

\textbf{E6 (ICON Velocity Tendencies, Figures~12--13, Table~IV):}
The velocity-tendencies module contains 25~loop nests that access 50~arrays.
After canonicalization, these collapse into 6~pattern classes with 7~layout-equivalent array groups.
Experiment~E6 evaluates a representative loop nest from each class individually and the full module end-to-end:
\begin{itemize}
\item \emph{Loopnest~1 (indirect stencil, complete vertical):} Vertical-first layout improves CPU bandwidth by $1.85\times$ on Zen4 (from 30\% to 55\% of STREAM peak) with a modest gain on Grace (from 44\% to 47\%).
On GPUs, both layouts achieve comparable bandwidth under exact R02B05 connectivity: 48--49\% on MI300A and 90--91\% on Hopper.
Under uniform random indirection, GPUs show large vertical-first gains: $3.3\times$ on MI300A and $2.3\times$ on Hopper.
The metric~$\mu$ is invariant to the indirection distribution under vertical-first layout on any platform (Table~IV).
\item \emph{Loopnests~2--6:} Representative nests from the remaining five pattern classes (direct stencil / partial vertical, direct stencil / complete vertical, indirect stencil + CFL clip / partial vertical, horizontal-only boundary, and vertical-only level reduction) are each implemented as a full C++/OpenMP + CUDA + HIP benchmark suite and an analytic cost-metric sweep.
Each suite sweeps V1--V4 dimension-order variants, 1D blocking $B \in \{8,16,32,64,128\}$, and (where vertical locality is meaningful) 2D tiling over $(T_X, T_Y) \in \{8,16,32,64\}^2$; GPU variants add a thread-block-size sweep.
The analytic cost metrics mirror the sweep axes and let us cross-check the empirical ordering against~$\mu$ and~$\Delta$.
\item \emph{Full module (GPU, cross-layout sweep):} Vertical-first yields $1.17\times$ on MI300A GPU and $1.07\times$ on Hopper GPU for Config~B (5 of every 6~calls); Config~A shows $1.06\times$ on MI300A and $0.97\times$ on Hopper.
The net per-timestep improvement is approximately 15\% on MI300A and 5\% on Hopper, weighted by the 5:1 call ratio.
The evaluation is driven by \texttt{run\_stage\{4,8\}\_permutations.py}, which sweeps a 95-configuration layout-permutation matrix (five dimension groups: \texttt{cv}, \texttt{ch}, \texttt{f}, \texttt{s}, \texttt{n}) per device, per loop-nest lowering stage.
On any platform, the aggregate improvement should be positive whenever the dominant configuration benefits from vertical-first layout.
\end{itemize}

\arttime
\begin{center}
\small
\begin{tabular}{lrr}
\toprule
Step & Daint & Beverin \\
\midrule
Setup                          &  15 &  15 \\
E1 Matrix Add.                 &  90 &  90 \\
E2 Conjugation                 & 210 & 210 \\
E3 Transpose                   & 180 & 180 \\
E4 GAS                         &  60 &  60 \\
E5 QE addusxx\_g               &  60 &  60 \\
E6 Loopnest 1 (indirect)       & 150 & 150 \\
E6 Loopnests 2--6              & 300 & 300 \\
E6 Full module sweep           & 240 & 240 \\
Analysis                       &  10 &  10 \\
\midrule
\textbf{Total}                 & \multicolumn{2}{c}{\textbf{$\sim$22~hr / cluster}} \\
\bottomrule
\end{tabular}
\end{center}
\noindent All times in minutes.
The E6 loopnests 2--6 line is an aggregate across five nests; each is a short-running sweep ($\sim$60~min) that can be scheduled independently.
The E6 full-module sweep assumes the three named configurations (\texttt{unpermuted}, \texttt{nlev\_first}, \texttt{index\_only}); a full 95-configuration sweep (\texttt{CONFIGS=full\_sweep}) is also supported but extends the full-module step to roughly 12~hr.

\artin

\artinpart{Hardware}

Two HPC clusters at the Swiss National Supercomputing Centre (CSCS):

\begin{itemize}
\item \textbf{Daint.Alps (Cluster~A):} Quad-GH200 nodes with 72-core Arm Neoverse~V2 Grace CPUs (512\,GB/s LPDDR5X) and NVIDIA H200 GPUs (4\,TB/s HBM3).
Measured STREAM Triad peaks: CPU 1.807\,TB/s, GPU 3.780\,TB/s.
\item \textbf{Beverin (Cluster~B):} Quad-MI300A APU nodes with 24 AMD Zen\,4 cores sharing 128\,GB unified HBM3.
Measured STREAM Triad peaks: CPU 1.161\,TB/s, GPU 4.294\,TB/s.
\end{itemize}

Both clusters are accessed via SLURM and require exclusive single-node allocations.
The results should generalize to any system with comparable cache-line sizes (64\,B on CPUs, 128\,B on GPUs), though absolute bandwidths will differ.

\artinpart{Software}

\begin{center}
\small
\begin{tabular}{lll}
\toprule
Package & Version & Notes \\
\midrule
GCC & 14.2 / 14.3 & via Spack \\
CUDA & 12.9 & Daint only \\
ROCm & 6.4.1 & Beverin only \\
HPTT & latest & E3 only \\
cuTENSOR & latest & E3, Daint only \\
hipTensor & latest & E3, Beverin only \\
DaCe & \texttt{yakup/dev} & E1--E5, E6 loopnests \\
     & \texttt{f2dace/staging} & E6 full module \\
Python & 3.10+ & numpy, scipy, matplotlib \\
libnuma & system & NUMA-aware allocation \\
\bottomrule
\end{tabular}
\end{center}

DaCe is available at \url{https://github.com/spcl/dace}.
Experiments E1--E5 and the E6 single-loopnest benchmarks use branch \texttt{yakup/dev}; E6's full-module evaluation uses branch \texttt{f2dace/staging} with layout transformations cherry-picked from \texttt{yakup/dev} (automated by \texttt{E6\_VelocityTendencies/full\_velocity\_tendencies/sc26\_layout/prepare.sh}).

Platform-specific setup and activation scripts are centralized under \texttt{Experiments/common/}:
\texttt{setup.sh} (first-time install), \texttt{activate.sh} (per-session), and the per-cluster environment scripts \texttt{setup\_daint.sh} / \texttt{setup\_beverin.sh}.
Every experiment's \texttt{run\_\{daint,beverin\}.sh} sources these in turn, so a single invocation of \texttt{common/setup.sh} at the repository root covers all experiments.

\artinpart{Datasets / Inputs}

Microbenchmarks (E1--E4) generate all input data synthetically at runtime.
The two production-code experiments require serialized datasets, each downloadable via a provided script:

\begin{itemize}
\item \textbf{E5 (Quantum ESPRESSO):} Serialized BaTiO$_3$ data from a QE \texttt{pw.x} HSE06 self-consistent calculation (10-atom unit cell).
Run \texttt{E5\_USXX/download\_data.sh} to fetch the data from ETH PolyBox.
%% TODO: insert final PolyBox share link once archived.
\item \textbf{E6 (ICON):} Serialized R02B05 grid connectivity tables (\texttt{p\_patch}), prognostic/diagnostic fields, and metric arrays from a single ICON timestep (timestep~9).
Run \texttt{E6\_VelocityTendencies/loopnest\_1/download\_data.sh} (or the equivalent \texttt{scripts/download\_nproma20480\_data.sh} in the full-module folder) to fetch the data from ETH PolyBox.
%% TODO: insert final PolyBox share link once archived.
\end{itemize}

\artinpart{Installation and Deployment}

Each experiment folder contains two self-contained scripts---\texttt{run\_daint.sh} and \texttt{run\_beverin.sh}---that handle the full pipeline: loading the platform environment, downloading data (if needed), compiling sources, executing benchmarks, and collecting results.
A typical script proceeds as follows:

\begin{enumerate}
\item Source the shared environment layer:
\begin{itemize}
\item \texttt{source ../../common/activate.sh} (activates the pyenv venv and pins DaCe to the branch specified by \texttt{\$DACE\_BRANCH}, defaulting to \texttt{yakup/dev}; \texttt{f2dace/staging} for E6 full module).
\item \texttt{source ../../common/setup\_\{daint,beverin\}.sh} (platform-specific Spack loads and OMP/NUMA pinning).
\end{itemize}
\item Download input data if not already present (E5 and E6 only).
\item Compile C++/CUDA/HIP benchmark sources (host tool\-chain selected by \texttt{CPU\_CXX}/\texttt{GPU\_CXX} from the platform env script).
\item Execute benchmarks with the configured OpenMP and NUMA settings (100~repetitions $+$ 5~warm-ups with inter-run cache flushes).
\item Collect CSV results into \texttt{results/\{beverin,daint\}/}.
\end{enumerate}

\noindent To reproduce a single experiment (e.g., E3 on Daint):
\begin{verbatim}
bash Experiments/common/setup.sh    # once per cluster
cd Experiments/E3_Transpose
sbatch run_daint.sh
# Results: results/daint/*.csv
python plot_paper.py                # produces Figure 9
\end{verbatim}

\artcomp

The artifact is organized as one folder per experiment.
Experiments E1--E5 each follow the same four-task workflow, encapsulated in their respective \texttt{run\_\{daint,beverin\}.sh} scripts:

\begin{description}
\item[$T_1$] \textbf{Setup:} source shared env, install DaCe, and download data (E5 only).
\item[$T_2$] \textbf{Build and execute:} compile sources and run benchmarks (100~repetitions after 5~warm-up iterations with cache flushes between repetitions).
\item[$T_3$] \textbf{Compute cost metrics} (E3 only): run \texttt{cost\_metrics} to compute $\mu$ and~$\Delta$ for each layout candidate.
\item[$T_4$] \textbf{Generate plots:} run \texttt{plot\_paper.py} to produce violin plots from the CSV results.
\end{description}

\noindent Dependencies: $T_1 \to T_2$; $T_1 \to T_3$; $\{T_2, T_3\} \to T_4$.
The run script handles $T_1$--$T_3$; $T_4$ is invoked manually afterwards.
Experiments E1--E5 are mutually independent and can run in parallel.

\textbf{E6 (ICON) workflow.}
Experiment~E6 follows a multi-stage pipeline that mirrors the paper's optimization methodology (Section~III-F):

\begin{description}
\item[$T_{6.1}$] \textbf{Download data:} fetch R02B05 connectivity tables and serialized fields via the per-loopnest \texttt{download\_data.sh} (or the full-module \texttt{scripts/download\_nproma20480\_data.sh}).

\item[$T_{6.2}$] \textbf{Access analysis:} \texttt{access\_analysis/analyze\_access.py} (plus \texttt{select\_loopnests.py}, \texttt{jaccard.py}, \texttt{coaccess.py}) canonicalizes the 25~loop nests into 6~pattern classes, groups the 50~arrays into 7~layout-equivalent groups, and computes $\mu$/$\Delta$ for each candidate layout per group.
Outputs the candidate configuration file \texttt{layout\_candidates.json}.

\item[$T_{6.3}$] \textbf{Per-loopnest evaluation:} each \texttt{loopnest\_\{1..6\}/} folder is a full bench\-mark suite (\texttt{bench\_cpu.cpp}, \texttt{bench\_gpu.cu}, \texttt{bench\_gpu\_hip.cpp}, \texttt{cost\_metrics.cpp}) that sweeps V1--V4 dimension orderings, 1D blocking, and (where meaningful) 2D tiling.
Each folder's \texttt{run\_\{daint,beverin\}.sh} compiles, runs, and writes CSVs to \texttt{results/\{daint,beverin\}/}; the per-folder \texttt{README.md} explains the Fortran reference nest and the mapping between the code's macros and the nest class.

\item[$T_{6.4}$] \textbf{Conflict resolution:} \texttt{conflict\_resolution/} reads the per-loopnest CSV rankings from $T_{6.3}$, weights them by call frequency (Config~A: 1/6, Config~B: 5/6), and emits \texttt{selected\_layouts.json}---the cross-product of locally-optimal layouts fed forward into $T_{6.5}$.
%% TODO: conflict-resolution implementation is currently scoped; only the README and input-schema contract are in place.

\item[$T_{6.5}$] \textbf{Full-module sweep:} \texttt{full\_velocity\_tendencies/} applies each selected layout from $T_{6.4}$ to the velocity-tendencies DaCe SDFG, drives the stage-by-stage lowering pipeline (stage1 $\to$ stage9), and benchmarks the terminal stage (stage4 on CPU, stage8 on GPU).
The driver is \texttt{run\_stage\{4,8\}\_permutations.py}; the clean submission scripts \texttt{run\_daint.sh} / \texttt{run\_beverin.sh} wire in \texttt{common/activate.sh} and \texttt{setup\_\{daint,beverin\}.sh}, set \texttt{DACE\_BRANCH=f2dace/staging}, and write outputs to \texttt{results/\{daint,beverin\}/stage\{4,8\}/}.
Outputs CSV results for Figure~13.

\item[$T_{6.6}$] \textbf{Analysis:} \texttt{plot\_paper.py} (top-level) and the per-folder plotting scripts under \texttt{full\_velocity\_tendencies/scripts/plotting/} produce Figures~12--13 and Table~IV from the combined results of $T_{6.3}$ and $T_{6.5}$.
\end{description}

\noindent Dependencies: $T_{6.1} \to T_{6.2} \to T_{6.3} \to T_{6.4} \to T_{6.5} \to T_{6.6}$.

\textbf{E6 sub-experiment structure.}
The ICON experiment spans two repositories and two DaCe branches.
Its folder is subdivided as follows:

\begin{itemize}
\item \texttt{access\_analysis/}: implements $T_{6.2}$.
Parses the annotated velocity-tendencies source (\texttt{velocity\_advection.f90}), canonicalizes the 25~loop nests into 6~pattern classes, groups the 50~arrays into 7~layout-equivalent groups, and emits \texttt{layout\_candidates.json}.

\item \texttt{loopnest\_\{1..6\}/}: one subfolder per pattern class, implementing $T_{6.3}$.
All six are complete C++/OpenMP + CUDA + HIP benchmark suites with matching analytic cost metrics.
The pattern classes are: (1) indirect stencil, complete vertical; (2) direct stencil, partial vertical; (3) direct stencil, complete vertical; (4) indirect stencil + CFL clip, partial vertical; (5) horizontal-only boundary; (6) vertical-only level reduction.

\item \texttt{conflict\_resolution/}: implements $T_{6.4}$.
Reads per-loopnest CSV results, applies call-frequency weighting, resolves cross-nest conflicts, and outputs the selected global layout configurations consumed by $T_{6.5}$.

\item \texttt{full\_velocity\_tendencies/}: implements $T_{6.5}$.
Contains the snapshot of the velocity SDFGs (pre-lowering baselines, per-stage checkpoints for stages~1--9, and the terminal stage5/stage8 program SDFGs in \texttt{SDFGs/}), the layout-transformation research code (\texttt{sc26\_layout/}), the DaCe-side transformation library (\texttt{scripts/utils/}), the permutation drivers (\texttt{scripts/run\_stage\{4,8\}\_permutations.py}), and the clean cluster submission scripts.
Uses repository~R2 with DaCe branch \texttt{f2dace/staging}.
\end{itemize}

\textbf{Statistical methodology.}
Following Hoefler and Belli~[52], each configuration runs 100~repetitions after 5~warm-up iterations.
Between repetitions, an $8192^2$ Jacobi stencil (3~sweeps with buffer swap, shared via \texttt{common/jacobi\_flush.h}) flushes all cache levels.
We report medians and 95\% bootstrap confidence intervals of the median (BCa method, $n=10{,}000$ resamples).
Violin plots show the full empirical distribution of all 100~samples; \emph{no outlier trimming is applied}.
The kernel density estimator uses Scott's rule for bandwidth and is evaluated at 200~equally-spaced points---the same settings in every figure---so violin shapes can be compared panel-to-panel and CPU-to-GPU without a hidden smoothing term distorting the comparison.
All random draws (input fills, index permutations, synthetic connectivity) are deterministic: every RNG is seeded from the canonical \texttt{SC26\_SEED=42} via a shared XOR-shift implementation (\texttt{common/prng.h}).
All bandwidths are normalized to the measured STREAM Triad peak on each platform.

\artout

Each experiment produces:
\begin{itemize}
\item CSV files with columns: variant, layout, repetition, time (ms), and bandwidth (GB/s).
\item PDF violin plots matching the corresponding paper figure (Figures~4, 8--13).
\item LaTeX-formatted metric tables where applicable (Tables~III and~IV).
\end{itemize}

Results are stored in \texttt{results/\{beverin,daint\}/} within each experiment folder.
The \texttt{plot\_paper.py} script reads these CSVs and produces PDF figures that can be compared visually with the paper.
Expected trends---the relative ordering of layouts and gap magnitudes---should be reproducible; absolute bandwidths will depend on the specific hardware and system load.

\end{document}
